% SOURCE_STATEMENT: "For equal-mass nonspinning ultrarelativistic black-hole collisions near critical impact parameter, does a unique critical solution exist, and can its high-energy behavior be described by an adiabatic, helical-Killing, or colliding-impulsive-wave approximation?"
% Preparation and final self-audit: CONTRIBUTING.md prompts C and D; source/audit artifacts retained in the local physics_statements_2026_10_11 batch.
% !TeX program = lualatex
\documentclass[11pt]{article}
\usepackage[margin=25mm]{geometry}
\usepackage{fontspec}
\setmainfont{Latin Modern Roman}
\usepackage{amsmath,amssymb,mathtools}
\usepackage{unicode-math}
\setmathfont{Latin Modern Math}
\usepackage{xurl,hyperref}
\hypersetup{colorlinks=true,urlcolor=blue}
\setcounter{secnumdepth}{0}
\setlength{\parindent}{0pt}
\setlength{\parskip}{5pt}
\setlength{\emergencystretch}{3em}
\newcommand{\sref}[2]{\hyperlink{source:#1:#2}{[S#2]}}
\newcommand{\cataloguescope}[1]{\paragraph{Recorded target.}#1}
\newcommand{\R}{\mathbb R}
\newcommand{\C}{\mathbb C}
\newcommand{\N}{\mathbb N}
\newcommand{\Z}{\mathbb Z}
\newcommand{\Q}{\mathbb Q}
\newcommand{\F}{\mathbb F}
\newcommand{\eps}{\varepsilon}
\begin{document}
% UnsolvedMath ID 30002147; source catalogue code OWR-12009-003
\setcounter{section}{30002146}
\section{Critical Ultrarelativistic Black-Hole Collision Solutions}\label{30002147}
{\small UnsolvedMath ID 30002147 · Mathematical Physics}
\subsection{Definitions and mathematical statement}

\paragraph{Shared notation.}$\mathbb N=\{1,2,\ldots\}$; the standard number fields have their usual meanings. All additional conventions are stated locally.


\paragraph{Definitions and conventions.}
The impact parameter is the transverse separation parameter of incoming equal-mass black holes in a scattering description. A critical impact parameter separates qualitatively different outcomes, such as scattering and merger. Nonspinning means zero initial intrinsic angular momentum; ultrarelativistic refers to large incoming boost or energy relative to rest mass. A helical Killing field combines time translation and rotation; an impulsive gravitational wave is a distributional spacetime idealization.

\paragraph{Target retained from the canonical source catalogue.}
\paragraph{Catalogue wording retained for traceability.}
For equal-mass nonspinning ultrarelativistic black-hole collisions near critical impact parameter, does a unique critical solution exist, and can its high-energy behavior be described by an adiabatic, helical-Killing, or colliding-impulsive-wave approximation?

\paragraph{Source scope and qualification.}
This is a numerical and analytical research programme. The source target leaves the limiting family, normalization and uniqueness notion to its collision model; no existence theorem for an unspecified class of initial data is asserted.
\subsection{Short English statement}
Is there a unique threshold geometry for high-energy black-hole collisions, and which proposed approximations describe it?
\subsection{Sources}
\begin{itemize}
\item[\hypertarget{source:30002147:1}{S1}] ULAM.AI and the UnsolvedMath Contributors, UnsolvedMath, record 30002147 (OWR-12009-003). Catalogue curation under CC BY 4.0; definitions and English exposition added locally.
\url{https://huggingface.co/datasets/ulamai/UnsolvedMath}.
\item[\hypertarget{source:30002147:2}{S2}] Mathematical Aspects of General Relativity (2012), Oberwolfach Reports, DOI 10.4171/owr/2012/37, pp. 2269--2333. Primary report accessed 2026-10-11; relevant PDF pages 31, 32, 33.
\url{https://doi.org/10.4171/owr/2012/37}.
\item[\hypertarget{source:30002147:3}{S3}] Klein, helical binary spacetimes (catalogue-listed context; not independently consulted in this preparation).
\url{https://arxiv.org/abs/gr-qc/0410095}.
\item[\hypertarget{source:30002147:4}{S4}] Healy--Ciarfella--Lousto, high-speed collisions (catalogue-listed context; not independently consulted in this preparation).
\url{https://doi.org/10.1103/PhysRevD.111.084077}.
\end{itemize}
\label{end:30002147}
\end{document}
